% ACCV 2026 — Rebuttal, Submission 230
%
% CHECK THE PAGE LIMIT before submitting. Numbers are reproduced from raw data by
% scripts/accv2026/rebuttal_*.py and checked by rebuttal_verify_numbers.py.
%
% Scope discipline: this answers what was asked. It does not volunteer internal
% issues no reviewer raised, and it does not include the anti-aliasing probe,
% which is inconclusive by construction (the models were fine-tuned on sharp
% frames, so filtered input is out-of-distribution) and would invite scrutiny
% without supporting a claim. Both are documented in the supplementary.
\documentclass[runningheads]{llncs}
\PassOptionsToPackage{table,dvipsnames}{xcolor}
\usepackage[review,year=2026,ID=230]{accv}
\usepackage{booktabs}
\usepackage{amsmath,amssymb}
\usepackage{xcolor}
\usepackage[hidelinks]{hyperref}

\setlength{\parskip}{2pt}
\setlength{\parindent}{0pt}
\newcommand{\TDS}{\textsc{TDS}}
\newcommand{\R}[1]{\textbf{R#1}}
\newcommand{\head}[1]{\vspace{2pt}\noindent\textbf{#1}\ }

\begin{document}
\pagestyle{empty}

\begin{center}{\large\bfseries Rebuttal --- Submission 230}\\[2pt]\end{center}

We thank the reviewers. \R{2} and \R{3} ask what our temporal protocol varies.
The question is well posed and we can now answer it precisely, with three new
experiments. The short answer is that the temporal axis measures \emph{evidence
availability} rather than sampling rate, that this is measurable and we quantify
it, and that the paper's conclusions hold under the corrected reading --- one of
them more strongly than before.

\head{What the protocol varies.}
Strided candidates are re-uniformised to each model's fixed frame budget $B$, so
at fixed coverage the selected frames span the same window at every stride, and
stride changes the input only once $\lceil\text{window}/s\rceil<B$. In seconds
between sampled frames (cov$=$100\%), UCF-101 reads $0.395/0.395/0.395/0.395/
0.617$ for $s{=}1..16$ and AUTSL $0.125/0.125/0.131/0.252/0.475$. We now state
this in Sec.~3.2 and report the temporal axis in seconds throughout.

\head{Three new experiments.}
\emph{(i) Sampling rate at fixed frame count.} Pinning $k$ to each model's budget
and varying only the window width $W$ changes the realised rate $k/W$ tenfold
while frame count is constant. Accuracy \emph{rises} with wider windows in 7 of 8
datasets (Spearman $+1.0$ in four; up to $+50.4$pp on AUTSL). Over the range this
protocol reaches, the binding constraint is how much of the action is observed.
\emph{(ii) Matched evidence.} Supplying every architecture the same $k$ distinct
frames at the same positions removes $B$ from the comparison. The architecture
ordering survives (Spearman $+0.857$, $p{=}0.007$, $n{=}8$ models) while the
budget correlation collapses from $+0.849$ ($p{=}0.008$) to $+0.231$ ($p{=}0.58$).
\emph{(iii) Multi-stride augmentation} (\R{1}'s request, below).

\head{What this changes, and what it does not.}
\TDS{} is unaffected: its dataset ranking is identical under curve-integral and
unclipped definitions (Spearman $1.000$) and swaps only adjacent pairs under
baseline- or chance-normalisation ($0.952$). The architecture finding stands but
is restated as evidence efficiency --- \emph{at an equal budget of distinct
frames, divided-attention and state-space designs deliver more accuracy per
frame} --- and we withdraw the $3$--$5\times$ magnitude, which reflected input
length rather than temporal aggregation. We also report a result the original
framing obscured: for a given dataset and frame budget there is a critical stride
at which a fixed-budget pipeline stops sampling and starts repeating, it varies
by an order of magnitude across datasets in seconds, and it is computable in
advance. We release that diagnostic.

\head{\R{3} (vdwp).}
\emph{Frame rate.} Correct. Measured from the files, three datasets differ from
the rates our duration analysis assumed (HMDB-51 30 not 25, AUTSL 30 not 25,
EPIC 50 not 60). At fixed nominal stride the realised interval varies $3.9\times$
across datasets, so we re-express every temporal axis in seconds.
\emph{Kinetics-400.} Added, with no training: all architectures are already
K400-pretrained and we verified label-order identity against each model's
\texttt{id2label}. It sits at the low-demand end, which sharpens the paper's
point --- the field's default recipes were tuned on the benchmark where sampling
matters least. \emph{Table 5.} FineGym added (Low $\leq50.4$pp, High $>71.2$pp,
33/32 classes), with the tier definition stated explicitly in the caption.
\emph{Table 7.} Your reading is right. Fine-tuning at the \emph{native}
resolution, where there is no shift to adapt to, already gains $+6$ to $+10.6$pp
for 7 of 8 models; the $+39.2$pp therefore decomposes into $\approx$32pp of
resolution adaptation and $\approx$7pp of additional training. We report both,
and add the without-fine-tuning row. \emph{Model ordering} is now consistent
across all figures and tables.

\head{\R{1} (XdvJ).}
\emph{Frames, span and interval versus stride:} answered above and now explicit
in Sec.~3.2. \emph{Multi-stride augmentation.} We ran it: two arms trained with
random $(\mathcal{C},s)$ exposure, against the published checkpoint. \textbf{Your
hypothesis is partly right, and more than we expected.} Exposure recovers
$28.4\%$ of the cliff for TimeSformer and $37.0\%$ for R3D-18, whose stride-16
accuracy rises $6.5\%\!\to\!31.9\%$. Two qualifications: 63--72\% of the cliff
persists, and the augmentation costs $3.3$/$9.9$pp of dense accuracy, so it
redistributes performance across operating points rather than removing the
limitation. At matched dense accuracy, the arm trained \emph{through} the
degraded input recovers $2.4$--$6.5\times$ more than one trained with uniform
resampling, which is what one expects if the residual is missing evidence rather
than a frequency effect. \emph{Convergence at 48/96px.} Recovered from 730
training logs: at 48px $52.8\%$ of runs still peak at the final epoch (16--27\%
elsewhere), so the budget does bind there, though the accuracy won in the last
three epochs is $1.38$pp, the smallest of any resolution. We state this as a
limitation. \emph{Sec.~3.3 PE interpolation} and the \emph{S19 SSv2 discussion}
are added.

\head{\R{2} (pxtb).}
\emph{Routing.} Both points correct. The rule thresholds $\max_c p(c|v)$, so
``entropy routing'' is a misnomer and is renamed. Choosing $\tau$ on a
calibration half and scoring on the held-out half (1000 stratified splits), the
SSv2 cell moves from $+0.21$pp to $-0.03$pp (95\% CI $[-1.07,+0.83]$): the
cascade \emph{matches} dense accuracy at $6.9$ frames rather than exceeding it.
We remove the ``halves temporal budgets without accuracy loss'' claim and
``Mitigating'' from the title, and present the cascade as a deployment
illustration. \emph{Statistics.} Correct that the 25 cells score the same clips.
Re-run as a within-subject design (clip as subject, each effect against its own
subject$\times$factor error), partial $\eta^2$ is coverage $0.292$, stride
$0.232$, interaction $0.070$; ``coverage is king at $2.2\times$'' becomes
$1.26\times$, while the per-model stride ordering is unchanged. On the
interaction: it \emph{is} reported (S13, FineGym 15.0\% of grid variance), though
we accept it was neither in the ANOVA model nor beyond two datasets, and we now
provide it for all 56 pairs, where it reaches $0.271$ (SlowFast/AUTSL).
\emph{\textsc{TDS}.} Each objection is testable and none changes the ranking:
$[\cdot]^+$ never fires, so removing it is numerically identical; the curve
integral gives Spearman $1.000$; baseline- and chance-normalised variants give
$0.952$. \emph{Configuration count} will be reported as the exact number of
distinct evaluations per protocol. \emph{Hardware scaling.} Correct --- the
device factors were estimates and are removed.

\end{document}
